% results_rq3

\section{RQ3: Travel and Competitive Edge}
\label{sec:rq3}

A third framing common in mobile-app HCI research is that differentiating
features win and retain users --- that a distinctive capability such as a
mosque finder or a travel prayer-shortening (\emph{qasr}) mode confers a
competitive edge. RQ3 tests this against what users actually cite as reasons
for loyalty, and against the sentiment record of apps that ship these two
travel-oriented features. Neither supports the differentiation story: what
earns loyalty in this ecosystem is the reliability of core functions, not
travel support specifically.

\subsection{What Users Cite Is Not Travel Support}

Across 4,696 preference citations in the corpus --- reviews explicitly stating
why the user chose or kept an app, in the register of ``that's why I use
this'' or ``only app that'' --- the aspects named most are adhan reminders
(374), Quran audio (293), prayer-time accuracy (290), qibla (155), and adhkar
(152). Mosque finder does not appear in the top eight cited aspects, and qasr
does not appear at all. Preference-citation rates are close to identical
whether an app ships a mosque finder or not: 1.38\% of reviews for the 12 apps
that have one, against 1.24\% for the 14 that do not. Whatever earns loyalty
in this ecosystem, it is not the travel feature set.

\subsection{The Mosque-Finder Coefficient Runs the Wrong Way, and the Reason Matters}

A cluster-robust fit on reviews mentioning mosque finder gives
\(\beta = -0.137\) (\(p = 0.003\), 21 app clusters, \(n = 3{,}010\)): apps that
\emph{have} a mosque finder attract more negative mentions of it, not fewer
--- 27.0\% negative against 23.0\% in apps without the feature, with positive
mentions moving the other way (38.6\% against 43.6\%). Before reading anything
into that direction, we note that the \texttt{mosque\_finder} tagger runs at
27\% precision (95\% CI [10, 57]), so of the 3,010 mentions entering
this comparison, drawn from 3,043 tagged corpus-wide, roughly 830 are genuine
(interval [297, 1{,}721]);
we do not report either figure unqualified anywhere in this paper. The noise falls on both arms of the
comparison rather than one, so the \emph{direction} of the gap is more
trustworthy than its exact \emph{size}.

The direction should not be read as the feature harming apps that ship it. The
two groups of mentions are not comparable speech acts. Of the 14 applications
without a mosque finder, 11 draw any mention of one, and there mentions are
overwhelmingly \emph{requests} --- 243 demand signals across 9 apps --- which
phrased as ``please add nearby mosques'' read neutral-to-positive. Of the 12
that have the feature, 10 draw mentions, and there they are \emph{usage
reports}; a usage report of a feature that fails to work reads negative. The
21 applications drawing mentions at all are the clusters in the model above. A user can only complain about a feature
that exists to complain about. The gap therefore measures the difference
between wishing for a feature and using one, not the cost of having shipped
it.

This coefficient is also unstable across specifications, which we take as
further reason not to over-read its direction: the mixed-effects fit on the
same contrast returns \(+0.05\) with an undefined (NaN) standard error, and
so cannot arbitrate between the two estimates. We treat the cluster-robust
estimate as the more trustworthy of the two, but rest no directional claim
about competitive advantage on a coefficient that reverses sign between
specifications, one of which failed to produce an error estimate at all.

\subsection{Qasr Travel Concession: Demand Evidence, Not a Causal Test}

Exactly one application of the 26, \emph{Muslim Bangla Quran Hadith Dua},
implements qasr (the Islamic travel concession shortening obligatory
prayers). With \(N = 1\), there is no between-app comparison available for
this feature, and we treat the qasr result as demand evidence rather than as
a test of any causal claim. Demand is real, if modest: 21 signals in total
--- 10 requests, 5 preference citations, 3 explicit statements of absence,
and 3 churn statements --- of which 20 come from 5 of the 25 apps that lack
the feature. The requests are specific and practical:

\begin{quote}
\itshape ``My second suggestion is that prayer times automatically adjust when
I'm traveling because man forgets and my prayers be off.''
\end{quote}
\hspace*{1em}--- Muslim Pro: Quran Athan Prayer, 4 stars

\begin{quote}
\itshape ``please add namaz qasar for all 5 prayers \ldots\ while anyone is out
of city or in traveling''
\end{quote}
\hspace*{1em}--- Prayer Times - Qibla \& Namaz, 5 stars

\begin{quote}
\itshape ``I wish it was installed prior my hajj trip''
\end{quote}
\hspace*{1em}--- Athan: Prayer Times \& Al Quran, 5 stars

A measurement caveat applies to the broader \texttt{qasr\_travel} aspect tag,
distinct from these 21 hand-read demand signals. That tag matches 637
mentions across 18 apps and runs 76.7\% positive against only 7.1\% negative
--- implausible for a capability 25 of the 26 apps lack --- and scored 0 of 2
on its gold-standard rows. Only 3.3\% of its matches are demand signals. The
lexicon is evidently catching travel and pilgrimage discourse in general
(hajj, umrah, journeys) rather than qasr functionality specifically. We
therefore never report 637 as a measure of qasr discussion volume, and rest
the qasr argument entirely on the 21 demand signals, which were read directly
rather than inferred from the lexicon match.

\subsection{Dominant Travel Failure Mode}

Where travel-related failures do appear, the recurring pattern is that prayer
times fail to adjust automatically when a user's location changes, and the
user discovers this only after acting on a stale time:

\begin{quote}
\itshape ``The only thing that I don't like about this app is it doesn't seem
to automatically change your location when traveling until you physically
refresh the prayer times page. Unfortunately today I accidentally broke my
fast 1 minute early b/c the location was still set to an area almost 1 hour
away, pls fix this inshaAllah.''
\end{quote}
\hspace*{1em}--- Athan: Prayer Times \& Al Quran, 4 stars

\subsection{Summary}

Neither travel-oriented feature confers a competitive edge. Mosque finder is
not among the top eight aspects users cite for loyalty, and its sentiment gap
reflects the difference between requesting and using a feature, not a cost of
shipping it. Qasr has essentially no market presence and rests on demand
evidence alone. What users do cite --- adhan reminders, Quran audio,
prayer-time accuracy --- points back to core functions, and those are
themselves among RQ6's most-broken promises (\S\ref{sec:rq6}). Competitive
advantage here looks like doing the basics reliably.
